\subsection{基于硬件 PMU 的 perf}\label{sec:perf}

要判断一段推理究竟慢在哪里，仅靠墙钟时间是不够的。墙钟只告诉我们一步花了多久，却说不清这段时间
里处理器到底在算还是在等。每一颗现代处理器内部都埋着一组性能监控部件（PMU，处理器内部的事件计数
器），它在不打扰程序运行的前提下，按硬件事件累加计数，比如经过的时钟周期、退休的指令条数、各级缓存
的访问与缺失、分支的执行与预测失败。RK3588 的 A55 与 A76 两种核都实现了 ARMv8 体系的第三代
监控规范 PMUv3，对外是一个专用的 64 位周期计数器 \texttt{PMCCNTR\_EL0} 与若干个可编程的 32 位
事件计数器 \texttt{PMEVCNTRn\_EL0}，每个可编程计数器配一个事件选择寄存器 \texttt{PMEVTYPERn\_EL0}
来指定它数哪一种事件，整组计数器再由 \texttt{PMCR\_EL0} 统一开关。Linux 上的 perf 工具正是读这些
计数器来给一段程序画像，把它的执行落到周期与指令这一层去解释。

Starry 此前只有软件与追踪类的事件，没有触达这组硬件计数器的通路。\textbf{本队为它补上了一个原生的 perf，
对齐 Linux perf 的语义，覆盖计数（\texttt{perf stat}）、采样（\texttt{perf record}）与符号化报告
（\texttt{perf report}）三种用法，使既有的 perf 二进制无需改动即可在 Starry 上对 RK3588 取数。}
这套能力本身是操作系统的一项通用设施，同时又是本项目优化所依赖的量具，第~\ref{sec:profiling}~节的
测量方法与第~\ref{sec:opt}~节的优化都建立在它读出的数字之上。它最初只能为单个任务在单核上计数，随后
扩展为对称多处理下的逐核采样，最终又针对 RK3588 的大小核异构结构，按集群分别配置各自的 PMU，并补上
必要的计数器复用与逐核分发。\textbf{这一系列工作以合并请求
\href{https://github.com/rcore-os/tgoskits/pull/1395}{\#1395} 提交至上游并已合并。}

\paragraph{把事件落到正确的 PMU 寄存器}
计数从识别事件开始。perf 递交的 \texttt{perf\_event\_attr} 既可能用 Linux 抽象的硬件事件号（如周期、
指令、缓存），也可能直接给出处理器原生的 ARM 事件号。前者须翻译为后者再写入 \texttt{PMEVTYPERn\_EL0}，
例如周期事件对应 ARM 事件号 \texttt{0x11}。周期事件若不带采样，则直接占用专用的 \texttt{PMCCNTR\_EL0}，
其余事件各取一个可编程计数器。一个看似简单却容易错的事件是分支指令计数，它在 A55 与 A76 上的最佳近似
并不相同，A55 倾向于用 \texttt{PC\_WRITE\_RETIRED}（\texttt{0x0C}），A76 则用 \texttt{BR\_RETIRED}
（\texttt{0x21}）。本实现据当前核所通告的事件支持位逐核裁决，而非全局写死一个号，否则在小核上数到的会是
另一回事。设备本身经 sysfs 暴露为 \texttt{armv8\_pmuv3\_0}，perf 据此发现这块 PMU 并选取事件。

\paragraph{从单核计数走到逐核分发}
\texttt{perf stat -a} 要对整机每一颗核分别取数，而发起调用的线程未必就跑在目标核上。计数器是处理器核
私有的，\texttt{PMEVCNTRn\_EL0} 只能在它所属的那颗核上配置与读取，跨核去碰别人的计数器毫无意义。为此每
颗核维护自己的一组计数器池，对某一目标核的配置、启动、读取、停止与释放，若发起者恰在该核便就地执行，否则
经一次同步的处理器间中断投递到目标核去做，发起线程阻塞到它返回，这与 Linux 用
\texttt{smp\_call\_function\_single} 把操作搬到目标核如出一辙。计数器的全局开关 \texttt{PMCR\_EL0.E}
也因此变成逐核各自置位。

\paragraph{计数器不够用时的轮换}
可编程计数器的数目有限，当一个任务同时打开的事件多过本核的计数器数目时，它们无法同时上硬件。此时各事件
轮流占用计数器，被注册到周期性的调度器节拍上，节拍每到一次就把当前任务的事件向前轮换一格，于是每个事件
都能在一段时间里被采到。轮换的代价是某个事件并非全程在数，实现因此分别记录每个事件处于使能状态的时长与
真正持有硬件计数器的时长，perf 拿到这两个时长后按使能比真实的比例把读数放大回去，这正是 Linux 计数器
复用的做法。采样路径则另有讲究，一个带采样周期的事件必定占用一个可编程计数器（即便是周期事件也不再走专用
计数器），并在计数溢出时触发中断，由溢出处理程序记一笔采样，处理程序运行在中断上下文里，不分配内存也不取
可睡眠的锁。

\paragraph{在两类核上分别取数}
RK3588 是大小核异构的，同一段推理放在 A55 与放在 A76 上，其周期与指令的关系并不一样，用一块全局 PMU
去笼统地数会把两类核的行为混在一起。本实现据 \texttt{MIDR\_EL1} 的实现者与部件号把每颗核归类到所属集群，
实现者为 Arm（\texttt{0x41}）、部件号 \texttt{0xD05} 即 Cortex-A55、\texttt{0xD0B} 即 Cortex-A76，
再把两个集群各自作为一块独立的 PMU 经 sysfs 暴露为 \texttt{armv8\_cortex\_a55} 与
\texttt{armv8\_cortex\_a76}。开在某一集群 PMU 上的事件只在该集群的核上配置与计数，于是同一段负载在
小核与大核上的画像可以分开取得并直接比较。

\paragraph{在本项目中读出的结论}
这套量具在本项目里有三种用法。其一，把周期与指令一起读出得到每周期指令数（IPC），再配上缓存事件，用以
判断一段负载是受算力约束还是受访存约束。其二，把流水线每一级的开销从墙钟挪到周期与指令这一尺度上，看清
时间究竟耗在何处。其三，在两个集群上分别读计数，量化同一段推理在 A55 与 A76 上的差异。借由它，本队得以
确认推理是受算力约束的，IPC 偏高而缓存缺失率极低，约 2.3\%，并非卡在访存，这一事实是把推理放置到大核的
直接依据，相关分析见第~\ref{sec:opt}~节。
